\subsection{8月}

\subsubsection{8月10日 $\sim$ 8月12日 周一$\sim$周三 —— SPI 硬件验证与模块开发}

\begin{itemize}
  \item \textbf{周一：组会 + SPI 硬件验证}：参加周一例会，并完成 SPI 接口的硬件验证工作，基于 FPGA 平台进行。
  \item \textbf{周二：clk\_gen 模块改写与验证}：完成 clk\_gen（时钟生成）模块的改写和验证，并将相关内容绘制到 PPT 中。同时完成本周周报的填写，发送给各位老师——\textbf{以后每周都要写周报}。
  \item \textbf{周三：ADC 行为级建模}：正在搭建 ADC 模块的行为级模型。
\end{itemize}

\begin{personalnote}
这周在硬件验证和模块开发上双线推进。SPI 验证依托 FPGA 平台跑通，clk\_gen 模块的改写和验证也顺利完成，还同步整理进了 PPT。周报从这周开始固定每周写，节奏要养成习惯。周三开始搭 ADC 的行为级模型，延续之前行为级建模的路线，后续应该会和校准模块的验证需求衔接上。
\end{personalnote}

\subsubsection{8月16日 周日 —— 想法：类 Verilog 语法糖的中间语言}

\begin{itemize}
  \item \textbf{背景痛点}：Scala 系 HDL（SpinalHDL、Chisel 等）生成的 Verilog 可读性差，对"翻译后"的 Verilog 不友好——结构不直观，难以直接修改或反查。
  \item \textbf{想法}：设计一种直观的、类似 Verilog 语法糖的中间语言（脚本），能够直观地展示出电路结构。
\end{itemize}

\begin{personalnote}
记录一个新想法：想设计一种"类 Verilog 语法糖"的中间语言。核心痛点在于 Scala 系 HDL（SpinalHDL、Chisel）生成的 Verilog 很不友好——生成结果可读性差，电路结构被抽象层掩盖，翻译/生成之后既难直接改、也难反查。这个中间语言的目标是：保留 Verilog 的直觉语法，但用语法糖把结构表达得更直观，让电路结构一眼可读。本质上是想在"高级 HDL 抽象"和"最终 Verilog"之间加一层可视化桥梁。后续值得想清楚它的定位：是纯前端表示，还是可逆的中间表示（能把生成的 Verilog 反编译回来）？
\end{personalnote}

\subsubsection{8月19日 周三 —— 新成员：喜乐蒂「缓存」到家}

\begin{itemize}
  \item \textbf{新成员}：买了一只喜乐蒂（Sheltie）狗狗，取名「缓存」。
  \item \textbf{置办}：购入大量宠物用品——狗粮、狗窝、玩具、牵引绳等，一次性配齐。
  \item \textbf{开销}：今天一共花了约 4000 元。
  \item \textbf{性格}：很乖，不乱叫；第一次坐摩托车就很配合，一路兜风回家。
\end{itemize}

\begin{personalnote}
今天家里多了一位新成员——一只喜乐蒂，取名叫「缓存」。搞硬件的人给狗起名都自带职业特色：cache 讲究的是命中率和低延迟，希望「缓存」在家里的表现也能又稳又快。小家伙确实争气：性格很乖，不吵不闹不乱叫，第一次坐摩托车就稳稳当当的，一路兜风回了家，适应性相当好。接下来要教的功课是定点上厕所，得花点耐心训练。不过迎接它的代价不小，一天下来狗粮、狗窝、玩具、牵引绳买了一大堆，钱包直接瘦了一圈，花了差不多 4000 块😭。但想想这也是长期投资，看着它蹦蹦跳跳的样子，这笔开销应该也算物有所值——接下来得努力搬砖攒狗粮钱了。
\end{personalnote}
